% Sources: README "Wind temperature and energy budget"; docs/wind_column_angles.md
% Code: ULXspectra, branch fix-results-integrity, commit 5f00569;
%       energy budget (Sect. wtemp:budget) and T_disc: ULXspectra-mc, branch mc-transfer-corrections, commit f546151

\section{Wind temperature and energy budget}
\label{sec:wtemp}

\subsection{Advection of photons by the wind}
\label{sec:wtemp:raw}

Photons released by the disc are carried outwards by the wind before they diffuse out of it (P07).
The photosphere above $R$ therefore shows the disc temperature of a smaller radius:
\begin{equation}
T_{\rm w,raw}(R) = T_{\rm disc}(\xeff\,\rin), \qquad \xeff = \frac{x}{1 + \taup v_{\rm w}/c} ,
\label{eq:Twraw}
\end{equation}
% code: wind_temperature_profile_array() (csrcs/wind_surface.c); README "Wind temperature"
where $T_{\rm disc}$ is the disc profile~\eqref{eq:Tdisc}, defined also beyond $\rout$, so that a hidden outer disc still sets the wind temperature.
$\xeff$ is never below $1$ or below the radius of the hottest disc ring.
% code: wind_temperature_profile_array() (x_eff_min = max(x_launch, x_Tmax)); README "Wind temperature"
Here $\taup$ is the vertical optical depth of the wind at $R$ and $v_{\rm w}$ its mean speed, so that $\taup v_{\rm w}/c$ is the ratio of the time of photon diffusion to the time in which the wind shifts them, $t_{\rm dif}/t_{\rm wind}$, for photons moving out of the disc in the vertical direction.
Eq.~\eqref{eq:Twraw} is a bridging formula, which reduces to the prescription of P07 at radii far enough from the centre.
% source: old manuscript (SEA_spectrum.zip, main.tex, "Effective temperature at the wind last-scattering surface"; Fig. photon_adv.png)
% VERIFY: the P07 equation it reduces to, and in which limit.
Each streamline crossing $R$ contributes with its own launch angle:
\begin{equation}
\taup = \frac{\kT}{4\pi R\,v_{\rm w}} \int_{\rin}^{\min(R, \Rsph)} \frac{{\rm d}\Mdot_{\rm w}(R_0)}{\sin\theta_0} ,
\label{eq:tauperp}
\end{equation}
% code: tau_perpendicular_simple_scalar(), wind_column_table_init() (csrcs/wind_surface.c); README "Wind temperature"; PLAN C14
with $v_{\rm w} = \beta_{\rm w}\sqrt{GM/\min(R, \Rsph)}$ (P07, eq.~27).
% code: v_wind_scalar() (csrcs/wind_struct.c)
The speed cancels in the advection ratio,
\begin{equation}
\frac{\taup v_{\rm w}}{c} = \frac{\kT}{4\pi R\,c} \int \frac{{\rm d}\Mdot_{\rm w}}{\sin\theta_0}
= \frac{\rg}{\etaEdd R} \int \frac{{\rm d}\mdot_{\rm w}}{\sin\theta_0} ,
\label{eq:advratio}
\end{equation}
% code: docs/wind_column_angles.md "Speed in the advection ratio (C17)"; PLAN C17
% derived: kappa_T Mdot_Edd/(4 pi c) = G M/(eta_Edd c^2) by eqs. (LEdd), (MEdd)
so the wind temperature depends on neither $\beta_{\rm w}$ nor the velocity law.

Eqs.~\eqref{eq:tauperp} and \eqref{eq:advratio} follow from the streamline density~\eqref{eq:rho}.
Streamlines do not cross (Eq.~\ref{eq:dzdx0}), so the points of the vertical line at $R$ inside the wind correspond one to one to the launch radii $\rin \le R_0 \le \min(R, \Rsph)$, with ${\rm d}z = |\partial z/\partial R_0|\,{\rm d}R_0$.
Since $|J| = \sin\theta_0\,|\partial z/\partial R_0|$, Eqs.~\eqref{eq:Sigmadotw} and \eqref{eq:rho} give
\begin{equation}
\kT\rho v\,{\rm d}z = \frac{\kT}{4\pi R}\,\frac{{\rm d}\Mdot_{\rm w}(R_0)}{\sin\theta_0} .
\label{eq:column_element}
\end{equation}
% derived (Claude) from Eqs. (Sigmadotw), (rho); generalises the old manuscript's derivation of tau_perp for
%       parallel streamlines (SEA_spectrum.zip, main.tex, "old Density in the wind")
For a speed $v = v_{\rm w}$ constant along the vertical, the integral over $z$ is Eq.~\eqref{eq:tauperp}.
For any velocity law, $\rho v$ is the mass flux along the streamlines, and $(1/c)\int\kT\rho v\,{\rm d}z$ equals the ratio~\eqref{eq:advratio}.
The advection ratio is therefore the one of the same density from which the photosphere is built.
In the code the two are computed independently, and at the NGC~7793 best fits they agree to 0.15 per cent.
% numbers: docs/wind_column_angles.md, table in "Speed in the advection ratio (C17)" (0.9985-1.0000), script docs/scripts/check_advection_ratio.py
% AUTHOR: give the commit for this check (docs say C17; CHANGELOG entry 2026-10-28).
Earlier versions of the code assigned all the wind at $R$ the angle $\theta_{\rm c}$; Eq.~\eqref{eq:tauperp} lowers $\taup$ by up to 20 per cent beyond $0.1\Rsph$ and changes the spectra at the NGC~7793 best fits by at most 0.056~dex.
% numbers: CHANGELOG 2026-10-28 "Streamline angle in the wind column", docs/wind_column_angles.md
% AUTHOR: keep this history in the paper, or drop it?
% OUTLINE: figure wind_column_angles.png (docs/scripts/compare_wind_column_angles.py, commit b6a9964).

\subsection{Energy budget}
\label{sec:wtemp:budget}

The wind radiates what the disc under its photosphere radiates:
\begin{equation}
L_{\rm wind} = L_{\rm disc,hidden} = 2\sum_{\rm hidden} \sigma T_{\rm disc}^4\,{\rm d}S_{\rm disc} ,
\label{eq:Lwind}
\end{equation}
% code: ULXspectra-mc f546151, bolometric_wind_power_from_surface() (L_disc_under_wind),
%       wind_temperature_correction_factor() (csrcs/wind_emission.c); README "Wind temperature and energy budget"
where the sum runs over the disc rings with $\rin \le R < \rout$ under the wind photosphere, on the same grid as the rest of the surface.
The temperatures~\eqref{eq:Tdisc} already exclude the wind toll, so $L_{\rm wind}$ is $(1 - \epsilon_{\rm w})$ of the release under the wind inside $\Rsph$ plus all of it beyond.
The disc rings that the optically thick wind does not cover keep their temperature~\eqref{eq:Tdisc}, also where they are seen through a transparent inner wind.
% source: README "Wind temperature and energy budget" ("Energy budget", "Directly seen disc")
% (Before 2026-11-03: L_wind = (1 - eps_w)(L_disc - L_disc,seen) with the analytic L_disc; CHANGELOG.)
The outer disc radius $\rout$ is set by the binary system and does not depend on the photosphere, whose extent depends on the mass loss in the wind.
A photosphere that reaches beyond $\rout$ has no disc under it there, and a disc that reaches beyond the photosphere is seen directly.
% source: the author 2026-11-09; code: csrcs/wind_emission.c (L_disc_under_wind: x_mid >= 1 and < x_out_disc,
%         "r_out, the disc's outer edge in the binary (user, 2026-11-03)"), ULXspectra-mc f546151; README "Energy budget"
%         (CHANGELOG 2026-11-03 "including rings beyond r_out" is superseded)
The budget is global: one factor scales the raw temperatures of all wind rings,
\begin{equation}
\Tw = f_{\rm cor}^{1/4}\,T_{\rm w,raw}, \qquad f_{\rm cor} = \frac{L_{\rm wind}}{L_{\rm wind,raw}} ,
\label{eq:fcor}
\end{equation}
% code: wind_temperature_correction_factor() (f_cor) (csrcs/wind_emission.c); README "Wind temperature and energy budget"
where $L_{\rm wind,raw}$ is the luminosity of both sides of the wind photosphere at the temperatures~\eqref{eq:Twraw},
\begin{equation}
L_{\rm wind,raw} = 2\int \sigma T_{\rm w,raw}^4\,2\pi R\,\sqrt{1 + \left(\frac{{\rm d}z_{\rm ph}}{{\rm d}R}\right)^2}\,{\rm d}R .
\label{eq:Lwindraw}
\end{equation}
% source: old manuscript (SEA_spectrum.zip, main.tex, "Luminosity of the last-scattering surface");
% code: dS_curve_one_side (csrcs/wind_surface.c), bolometric_wind_power_from_surface() (csrcs/wind_emission.c)
The integral runs along the photosphere including the outer wall, where it becomes $2\pi R_{\rm wall}\,{\rm d}z$.
% code: wind_model_api.c (include_outer_vertical_wall = 0: the wall is part of the surface curve)
$\Tw$ is the temperature of a blackbody with the energy of the wind ring; the bolometric luminosities and the transfer correction (Sect.~\ref{sec:spectrum:transfer}) use it, while the spectrum of the wind ring is a diluted blackbody at $T_{\rm w,raw}$ (Sect.~\ref{sec:spectrum}).
% source: README "Wind temperature and energy budget" ("Energies (L_bol, the transfer correction) use T_w in both cases")
Since the disc temperature steps at $\Rsph$, the wind temperature steps where $\xeff = \xsp$, that is at $\Rsph(1 + \taup v_{\rm w}/c)$, which is 1.2--2.2~$\Rsph$ at the old best fits.
% source: README "Wind temperature and energy budget" ("Where eps_w is taken"); ULXspectra-mc f546151
$f_{\rm cor}$ is not capped.
A value $f_{\rm cor} > 1$ means that the wind radiates more than its surface would at $T_{\rm w,raw}$; it is allowed and the model is flagged.
If $L_{\rm wind} = 0$ ($\epsilon_{\rm w} = 1$, or no disc ring under the wind), the wind is dark, $f_{\rm cor} = 0$, and the model is flagged as well.
% code: flags wind_f_cor_above_one, wind_budget_fallback (diagnostics and MCMC blobs); README "Wind temperature and energy budget";
%       CHANGELOG 2026-11-04 (cap removed), 2026-11-03 (dark wind instead of f_cor = 1)
% AUTHOR: is a flagged model rejected in the fits, or kept? (Sect. fitting)

The model does not determine the wind toll $\epsilon_{\rm w}$.
Energy conservation allows any value between 0 and 1; at $\epsilon_{\rm w} = 1$ the whole release inside $\Rsph$ is paid as the toll, and no photons are left from there.
% source: the author 2026-11-09
$\epsilon_{\rm w}$ is therefore a free parameter, and in the fits it acts mainly on the brightness of the wind: where $\epsilon_{\rm w}$ is much larger than the advection ratio $\taup v_{\rm w}/c$, it works as a free dimmer of the wind.
A lower bound on the toll from the energy the wind needs is discussed in Sect.~\ref{sec:disc:epsw}.
% source: ULXspectra-mc PLAN.md C21 (check of eps_w against tau_perp v_w/c, 2026-11-04: 8 of 25 fits with eps_w free,
%         f_cor <~ 0.05 there; "eps_w is treated as phenomenological"); fits from before the step at R_sph
% AUTHOR: a lower bound on eps_w from the energy the wind carries out (kinetic + enthalpy flux through the photosphere,
%         diagnostic wind_energy_loss_through_green_surface(), Discussion outline) - compare in the Discussion?
